\setcounter{chapter}{6}
\chapter{Quicksort}
%Section 7.1------------------------------------------------------------------------------------------------------------
\section{Description of quicksort}
\subsection*{Note}
\begin{enumerate}
%
\item An alternative to the loop invariant on page 184 (for the \kw{for} loop of lines 3--6 in $\proc{Partition}$) is
\begin{loopinv}
At the beginning of each iteration of the loop of lines 3--6, the following hold:
\begin{enumloopinv}
%%
\item The subarray $A[p \subarr i]$ contains exactly the elements originally in $A[p \subarr i]$ (prior to execution of the loop) that are $\leq x$.
%%
\item The subarray $A[i + 1 \subarr j - 1]$ contains exactly the elements originally in $A[i + 1 \subarr j - 1]$ (prior to execution of the loop) that are $> x$.
%%
\item $A[r] = x$.
%%
\end{enumloopinv}
\end{loopinv}
%
\end{enumerate}
\subsection*{Exercises}
\begin{enumerate}[\thesection-1]
%
\setcounter{enumi}{1}
%
\item When all elements in the subarray $A[p \subarr r]$ have the same value, the returned value $i + 1$ for $q$ is $r$. The procedure modifies $\proc{Partition}$ so that it returns $q = \floor{(p + r) / 2}$ when all elements in the subarray $A[p \subarr r]$ have the same value.
\begin{codebox}
\Procname{$\proc{Partition}'(A, p, r)$}
\li $x \gets A[r]$
\li $i \gets p - 1$
\li $\id{all-the-same} \gets \const{true}$
\li \For $j \gets p$ \To $r - 1$
\li \Do
         \If   $A[j] \neq x$
\li      \Then
               $\id{all-the-same} \gets \const{false}$
         \End
\li      \If   $A[j] \leq x$
\li      \Then
               $i \gets i + 1$
\li            exchange $A[i]$ with $A[j]$
         \End
    \End
\li exchange $A[i + 1]$ with $A[r]$
\li \If $\id{all-the-same} \isequal \const{true}$
\li \Then \Return $\floor{(p + r) / 2}$
\li \Else \Return $i + 1$
    \End
\end{codebox}
%
\item The loop of lines 3--6 runs in exactly $n - 1$ iterations.
%
\item Change $\leq$ to $\geq$ in the \kw{if} test in line 4.
%
\end{enumerate}
%End of Section 7.1-----------------------------------------------------------------------------------------------------


%Section 7.2------------------------------------------------------------------------------------------------------------
\section{Performance of quicksort}
\subsection*{Exercises}
\begin{enumerate}[\thesection-1]
%
\setcounter{enumi}{3}
%
\item As noted in the last paragraph of the subsection ``Worst-case partitioning,`` when the input array is almost sorted, $\proc{Insertion-Sort}$ runs in linear time while $\proc{Quicksort}$ runs in quadratic time.
%
\setcounter{enumi}{5}
%
\item Such a split occurs when the last element used as the pivot is not among the highest or lowest $\alpha$ portion of all elements in the array.
%
\end{enumerate}
%End of Section 7.2-----------------------------------------------------------------------------------------------------


%Section 7.3------------------------------------------------------------------------------------------------------------
\section{A randomized version of quicksort}
\subsection*{Exercises}
\begin{enumerate}[\thesection-1]
%
\item The worst-case running time of a randomized algorithm is asymptotically equal to that of a deterministic (i.e., not randomized) algorithm, assuming that the time taken by randomization is dominated by the worst-case running time of the deterministic algorithm.
%
\item Assume that the input array has length $n$. Note that the best-case and worst-case behaviors of $\proc{Randomized-Quicksort}$ are asymptotically the same as those for the deterministic version, $\proc{Quicksort}$. As noted in Section 7.2 of the text, the worst-case scenario happens when the maximally unbalanced split occurs at every level of the recursion associated with $\proc{Quicksort}$, thus a total of $n - 1 = \bigtheta{n}$ calls are made to $\proc{Random}$. On the other hand, the best-case scenario is when (almost) perfectly balanced split occurs at every level of the recursion, which gives rise to a roughly complete binary tree for its recursion tree; there are approximately $n / 2$ internal nodes, resulting in a total of $\bigtheta{n}$ calls to $\proc{Random}$.
%
\end{enumerate}
%End of Section 7.3-----------------------------------------------------------------------------------------------------


%Section 7.4------------------------------------------------------------------------------------------------------------
\section{Analysis of quicksort}
\subsection*{Note}
\begin{enumerate}
%
\item Line 7 of the proof of Lemma 7.3~on page 196: Change $\proc{Randomized-Select}$ to $\proc{Randomized-Partition}$.
%
\end{enumerate}
\subsection*{Exercises}
\begin{enumerate}[\thesection-1]
%
\setcounter{enumi}{1}
%
\item (Sketch). Let $T(n)$ be the best-case running time of quicksort, then
\[
T(n) = \min\setm{T(q) + T(n - q - 1)}{0 \leq q \leq n - 1} + \bigtheta{n}.
\]
Assume $T(n) \geq c n \lg n$ for some suitable $c > 0$, we have
\[
T(n) \geq c \multp \min\setm{q \lg q + (n - 1 - q) \lg (n - 1 - q)}{0 \leq q \leq n - 1}.
\]
Consider the function $f(x) = x \lg x + (n - 1 - x) \lg (n - 1 - x)$ whose domain consists of \emph{real numbers} in $(0, n - 1)$; note that $f(n - 1 - x) = f(x)$ for $x \in (0, n - 1)$. It follows that $f'(x) < 0$ for $x \in (0, (n - 1)/2)$ and $f'(x) > 0$ for $x \in ((n - 1)/2, n - 1)$; moreover $f'((n - 1)/2) = 0$.
%
\setcounter{enumi}{3}
%
\item We can derive the lower bound $\bigomega{n \lg n}$ for the expected running time of $\proc{Randomized-Quicksort}$  analogously to the textbook, by adopting the following changes:
\begin{enumerate}[(1)]
%%
\item Change $\bigoh{n + X}$ to $\bigomega{n + X}$ in Lemma~7.1: Due to the condition in line 1 of $\proc{Quicksort}$, for all $1 \leq i < n$ we have that at least one between $A[i]$ and $A[i + 1]$ is selected as a pivot in an execution of $\proc{Partition}$. Thus, there can be at least $\floor{n / 2}$ calls to $\proc{Partition}$ over the entire execution of the quicksort algorithm, and hence there are at least $2 \multp \floor{n / 2}$ recursive calls to $\proc{Quicksort}$ itself.
%%
\item Change $\bigoh{n \lg n}$ to $\bigomega{n \lg n}$ in Theorem~7.4: Use the following in place of the inequality at the bottom of page~197 in textbook:
\[
\ds{\sum^{n - 1}_{i = 1} \sum^{n - i}_{k = 1} \frac{2}{k + 1} \geq \sum^{\floor{n / 2}}_{i = 1} \sum^{\ceil{n / 2}}_{k = 1} \frac{2}{k + 1} = \sum^{\floor{n / 2}}_{i = 1} \parenf{-2 + \sum^{\ceil{n / 2}}_{k = 0} \frac{2}{k + 1}}}
\]
for sufficiently large $n$.
%%
\end{enumerate}
%
\item (INCOMPLETE)
%
\item There are ${n \choose 3}$ possible combinations of choosing three elements randomly from the input array. For ease of analysis, it is convenient to assume that all of the elements are distinct; more precisely, let $z_1, z_2, \etl, z_n$ be the elements in the input array with $z_1 < z_2 < \etc < z_n$. The desired probability can be evaluated by noting the two mutually exclusive events below:
\begin{enumerate}[(1)]
%%
\item The median of the three randomly chosen elements is $z_k$ for some $2 \leq k \leq \alpha n$. Then the smallest and the largest elements among the three are then $z_i$ for some $1 \leq i < k$ and $z_j$ for some $k < j \leq n$, respectively.
%%
\item This median of the three randomly chosen elements is $z_k$ for some $(1 - \alpha)n + 1 \leq k \leq n - 1$. The smallest and the largest elements among the three are $z_i$ for some $1 \leq i < k$ and $z_j$ for some $k < j \leq n$, respectively.
%%
\end{enumerate}
Note that the above two events have the same probability.
%
\end{enumerate}
%End of Section 7.4-----------------------------------------------------------------------------------------------------


%Problems of Chapter 7--------------------------------------------------------------------------------------------------
\section*{Problems}
\begin{enumerate}[\thechapter-1]
%
\item
\begin{enumerate}[\it a.]
%%
\setcounter{enumii}{1}
%%
\item Assume that the elements in the subarray $A[p \subarr r]$ are all equal, and let $n \defas r - p + 1$, i.e., the number of elements in this subarray. Then, $\proc{Hoare-Partition}$ differs from $\proc{Partition}$ in that the former returns $\floor{(p + r)/2}$ (with $j = (p + r - 1)/2 = i - 1$ if $n$ is even and $j = (p + r)/2 = i$ if $n$ is odd) while the latter returns $r$. Thus, $\proc{Hoare-Partition}$ performs better than $\proc{Partition}$ in this situation in that the former creates a perfectly balanced partition of $A[p \subarr r]$ while the latter creates a most unbalanced partition.
%%
\end{enumerate}
%
\end{enumerate}